\begin{textsample}{POS Dim 2 – pb – Score 17.00 – pb/pb\_em\_114.txt}  \label{ex:f2_pos_280}
4.2.4.2 Formação para o Mundo do Trabalho O processo de formação de um \textbf{profissional} deve envolver mais do que as competências específicas do \textbf{curso} \textbf{técnico}.
Faz -se necessário trabalhar com competências gerais para o mundo do trabalho, que envolve desde as novas tecnologias até as competências socioemocionais.
A Formação Básica para o Mundo do Trabalho contempla os 4 eixos \textbf{estruturantes}, a partir da \textbf{oferta} das três disciplinas empreendedoras : Inovação Social e Científica, Intervenção Comunitária, Empresa Pedagógica, Educação Tecnológica e Midiática e Higiene e Segurança do Trabalho.
Para implementar a Formação Básica para o Mundo do Trabalho, é necessário desenvolver formação docente continuada, considerando as especificidades de cada modalidade da \textbf{oferta}, as ferramentas necessárias, o planejamento e a avaliação.
Além disso, é importante que a escola disponha de laboratório de informática e laboratório maker, visto que os estudantes irão explorar diversos contextos e lidar com várias problemáticas, em que eles precisarão de um espaço de prática com ferramentas e recursos básicos para gerar uma proposta de solução. 4.2.4.2.1 Disciplinas empreendedoras ( Eixos \textbf{Estruturantes} ) I.
História da criação das disciplinas empreendedoras A partir da \textbf{implantação}, em 2016, das oito Escolas Cidadãs Integrais ( ECIs ), das quais três eram \textbf{técnicas} ( ECITs ), a Rede Estadual de Educação da Paraíba iniciou a \textbf{oferta} de \textbf{cursos} \textbf{técnicos} integrados ao Ensino Médio, compostos por disciplinas propedêuticas e \textbf{técnicas}, além de uma parte diversificada com opções de escolha pelo estudante de disciplinas \textbf{eletivas} e clubes juvenis.
No entanto, durante o processo de \textbf{implantação}, houve inquietação pela equipe pedagógica de \textbf{implantação} acerca das ECITs seguirem o mesmo modelo pedagógico inovador das ECIs na parte diversificada e, na parte \textbf{técnica}, se diferenciarem apenas pela \textbf{oferta} de componentes curriculares específicos.
Faltavam ações que mantivessem a escola \textbf{sintonizada} com desafios comunitários e empreendedores relacionados às atividades do mundo do trabalho, além de garantir um foco na parte diversificada capaz de promover o protagonismo \textbf{profissional} e social do estudante.
Dessa forma a equipe da Secretaria de Estado da Educação ( SEE ) optou por realizar um projeto piloto voltado para as ECITs, o qual articularia o \textbf{currículo} \textbf{técnico} com o propedêutico, com a preparação básica para o trabalho e com a parte diversificada do \textbf{currículo}, além de criar disciplinas \textbf{eletivas} voltadas ao protagonismo \textbf{profissional} e social.
Nesse processo, surgiu o Projeto Empreendedor no \textbf{currículo} das escolas \textbf{técnicas} e, junto a ele, as disciplinas empreendedoras : Inovação Social e Científica ( ISC ), Intervenção Comunitária ( IC ) e Empresa Pedagógica ( EP ).
O projeto, que é \textbf{ofertado} na \textbf{rede} \textbf{estadual} desde 2018, foi construído pelos professores, por meio de formações, com o objetivo de elaborar um \textbf{currículo} articulado pautado em competências e habilidades e \textbf{ofertar} aos participantes um arcabouço teórico conceitual que permitisse uma compreensão da prática cotidiana sobre esses conceitos.
Desse modo, o \textbf{currículo} se expressa, na sala de aula, sempre por meio de atividades para os estudantes desenvolverem.
As disciplinas integram as habilidades dos quatro eixos \textbf{estruturantes} propostos pela lei do Novo Ensino Médio e fortalecem a proposta de utilização de parte da \textbf{carga} \textbf{horária} da Preparação Básica para o Trabalho do \textbf{currículo} dos \textbf{cursos} \textbf{ofertados} nas ECITs, em uma ação interdisciplinar prática inovadora, de aprofundamento dos conhecimentos e aquisição de competências.
Atualmente as disciplinas empreendedoras são \textbf{ofertadas} em todas as Escolas Cidadãs Integrais \textbf{Técnicas} da Rede Estadual, visto que trabalham o processo de articulação curricular, ao abarcar diversas competências, temáticas e conceitos dos \textbf{cursos} \textbf{técnicos} e das Áreas de Conhecimento, bem como os 4 Eixos \textbf{Estruturantes} propostos pelo Novo Ensino Médio.
Diante disso, as disciplinas deverão estar dentro do Módulo de Formação Básica para o Trabalho, em todas as \textbf{ofertas} como parte do Itinerário de Educação Profissional e \textbf{Técnica}, seguindo as exigências da reformulação do Ensino Médio.
II.
Disciplinas empreendedoras As disciplinas empreendedoras dialogam com os eixos \textbf{estruturantes}.
Suas ementas estão disponíveis e as orientações de cada elemento da sequência didática também, para que as escolas de educação \textbf{profissional} públicas ou privadas possam trabalhar os eixos.
As metodologias criadas de forma coletiva visam proporcionar a aquisição de competências relacionadas ao mundo do trabalho.
Logo, as execuções dos projetos empreendedores não precisam necessariamente relacionar -se com o \textbf{curso} \textbf{técnico} específico oferecido pela escola, mas sim permitir que o estudante vivencie atividades práticas de campo, de tal forma que aplique os conhecimentos da parte propedêutica do \textbf{currículo} e compreenda o contexto onde está inserido, a cadeia produtiva de uma empresa, as relações de trabalho e interpessoais existentes, os serviços comunitários etc.
É importante frisar duas bases fortes no corpo das disciplinas empreendedoras : a ) Planejamento pautado em competências e habilidades : com o objetivo de elaborar o \textbf{currículo} articulado embasado em competências e habilidades, sempre por meio de atividades para os estudantes desenvolverem, como já argumentado anteriormente ; b ) Métodos de resolução de problemas : as habilidades para o protagonismo \textbf{profissional}, identificadas nas visitas realizadas às empresas, deixaram evidente que, para se tornar um adulto emancipado, o estudante precisa ter capacidade de resolver problemas das mais diversas naturezas ; lidar com desafios pessoais, como desenvolver autonomia, resiliência e capacidade de liderança, além de enfrentar a frustração e desafios \textbf{profissionais}, como ter boa comunicação, pensamento estratégico, visão sistêmica, flexibilidade, capacidade de análise de dados, raciocínio lógico, além de saber propor soluções, construir textos etc.
Com a proposta do novo ensino médio a partir da lei no 13.415/2017, em que o \textbf{itinerário} de educação \textbf{técnica} deve contemplar os eixos \textbf{estruturantes} previstos, as disciplinas vão integrar a formação básica para o trabalho, que irá ser contemplado em todas as propostas de educação \textbf{profissional}, sendo elas FIC, integral ou concomitante.
As disciplinas empreendedoras são : - Intervenção Comunitária : visa a uma mudança na comunidade que promova o bem-estar das pessoas por meio da aplicabilidade das competências da Base Comum Curricular e do \textbf{curso} \textbf{técnico}.
O Projeto de Intervenção Comunitária poderá ser desenvolvido em consonância com empresa parceira, com a Prefeitura, com o Estado ou somente com a escola. - Inovação Social e Científica : permite a participação dos estudantes em \textbf{oficinas} práticas, facilitadas pelos professores, cujo tema gerador é o desenvolvimento de tecnologias sociais em áreas como Ciência, Tecnologia, Engenharia e Matemática relativas ao \textbf{curso} \textbf{técnico} escolhido. - Empresa Pedagógica : visa à aquisição de competências e habilidades gerais para melhor desempenho do estudante no mundo do trabalho.
Bem como trabalhar o \textbf{empreendedorismo} pessoal, a partir da educação financeira, compreensão do desenvolvimento econômico e pessoal.
Pensar uma educação atual requer investir, experimentar e acreditar que é possível mudar mesmo dentro das limitações das escolas, sejam privadas ou públicas.
Precisa -se, portanto, de \textbf{profissionais} que acreditem que a escola deve abrir as portas para interagir com o entorno e que o estudante é capaz de construir a própria história, o seu projeto de vida ; e orientá -lo para isso.

% matched lemmas: carga, currículo, curso, eletivo, empreendedorismo, estadual, estruturante, horário, implantação, itinerário, oferta, ofertar, oficina, profissional, rede, sintonizar, técnico
\end{textsample}
